\subsection{Circulares Únicas}

{Circular Única de Casas de Bolsa}

En la Circular Única de Casas de Bolsa, los artículos relacionados con la administración del riesgo de mercado se encuentran principalmente en el \textbf{Título Quinto, Capítulo Tercero}, relativo a la \textbf{Administración Integral de Riesgos}. Dentro de esta sección, el \textbf{artículo 122} clasifica los riesgos que enfrentan las casas de bolsa y ubica al \textbf{riesgo de mercado} como un riesgo cuantificable discrecional, mientras que los \textbf{artículos 123 al 133} establecen la estructura general de gobierno, políticas, límites, metodologías, modelos, seguimiento y auditoría que deben sostener su administración. En particular, la definición general de administración integral de riesgos se entiende como el conjunto de objetivos, políticas, procedimientos y acciones para identificar, medir, vigilar, limitar, controlar, informar y revelar los distintos riesgos a que están expuestas las casas de bolsa.

De manera específica, los \textbf{artículos 138, 139, 140 y 141} contienen el núcleo normativo sobre administración del riesgo de mercado. El \textbf{artículo 138} dispone que, respecto de títulos para negociar, títulos disponibles para la venta, reportos, otras operaciones con valores e instrumentos derivados de negociación y de cobertura, las casas de bolsa deben analizar, evaluar y dar seguimiento a sus posiciones mediante \textbf{modelos de Valor en Riesgo (VaR)}, capaces de medir la pérdida potencial asociada a movimientos de precios, tasas de interés o tipos de cambio. Asimismo, este artículo exige asegurar la consistencia entre modelos de valuación, evaluar concentraciones, comparar las exposiciones estimadas con los resultados observados, contar con información histórica de factores de riesgo y calcular pérdidas potenciales bajo distintos escenarios, incluidos escenarios extremos.

Por su parte, el \textbf{artículo 139} regula la administración del riesgo de mercado para títulos conservados a vencimiento, derivados de cobertura de posiciones distintas a las del artículo 138 y demás posiciones sujetas a este riesgo. En este caso, la medición debe enfocarse en las variaciones de los ingresos financieros y del valor económico, utilizando modelos capaces de estimar la pérdida potencial derivada de movimientos en tasas de interés por moneda. Además, se exige consistencia entre modelos, comparación entre resultados estimados y observados, uso de información histórica y evaluación bajo escenarios extremos.

El \textbf{artículo 140} prevé que ciertos títulos clasificados como disponibles para la venta puedan administrarse bajo el régimen del artículo 139, en lugar del artículo 138, siempre que exista aprobación del comité de riesgos, se justifique que dichos títulos forman parte estructural del balance y se establezcan controles internos adecuados. A su vez, el \textbf{artículo 141} señala que, para reconocer el propósito único de cobertura de un instrumento financiero y sujetarlo al tratamiento de los artículos 138 o 139, debe demostrarse una relación inversa significativa entre el instrumento de cobertura y la posición cubierta, sustentada con evidencia estadística suficiente y con seguimiento continuo a la efectividad de la cobertura.

De forma complementaria, los \textbf{artículos 127 al 132} fortalecen la administración del riesgo de mercado al establecer que el comité de riesgos debe aprobar objetivos, lineamientos, políticas y límites globales y específicos de exposición a los distintos tipos de riesgo, así como mecanismos de acción correctiva y supuestos de excepción. Finalmente, el \textbf{artículo 144} exige la revelación cuantitativa de los riesgos asumidos por la casa de bolsa, incluyendo el \textbf{Valor en Riesgo de mercado}, la evaluación de variaciones en ingresos financieros y valor económico, así como valores promedio de exposición por tipo de riesgo, lo cual completa el marco normativo de medición, control y transparencia del riesgo de mercado.